%% ma: L11-B25-0002
%% lop: 11
%% chuong: 7
%% bai: 25
%% dang: 11.25.2
%% loai: TLN
%% muc: VD
%% dap_an: "1,09"
%% nguon: "(NBV)-ĐỀ SỐ 1-MỖI NGÀY 1 ĐỀ THI 2025.pdf (D01, MỖI NGÀY 1 ĐỀ - Đề số 1), Phần III Câu 1 (THPT Đào Duy Từ - Thanh Hóa 2025)"
%% kien_thuc: "Lớp 11 Bài 25 (góc nhị diện, góc giữa hai mặt phẳng), Bài 23-24 (đường thẳng vuông góc với mặt phẳng, hình chiếu vuông góc); hình chóp tứ giác đều"
%% ghi_chu: "Đề gốc có ảnh minh họa kim tự tháp (không cần cho lời giải nên không vẽ lại). Đáp án gốc 1,09 đúng"
%% trang_thai: cho-duyet
%% ly_do: "Dạng mới đề xuất 11.25.2 (góc giữa mặt bên và mặt đáy của hình chóp tứ giác đều, số liệu thực tế) chờ giáo viên duyệt"
\begin{ex}
Kim tự tháp Memphis tại bang Tennessee (Mỹ) có dạng hình chóp tứ giác đều với chiều cao $98$ m và cạnh đáy $180$ m. Gọi $\alpha$ là góc nhị diện tạo bởi mặt bên và mặt đáy của kim tự tháp. Tính $\tan\alpha$. (Làm tròn kết quả đến hàng phần trăm.)
\shortans{1,09}
\loigiai{Gọi $S.ABCD$ là hình chóp, $O$ là tâm đáy, $M$ là trung điểm cạnh $AB$. Khi đó $OM\perp AB$ và $SM\perp AB$ nên $\alpha=\widehat{SMO}$.
Có $SO=98$, $OM=\dfrac{180}{2}=90$ nên $\tan\alpha=\dfrac{SO}{OM}=\dfrac{98}{90}=\dfrac{49}{45}\approx1{,}09$.}
\end{ex}
